\clearpage
\fancyhead[L]{\small CONNAISSANCE DES TEMS FÜR 1818.}
\hypertarget{native-600-lower}{}
\newcommand{\ObservedJupiterErrorMark}{\hbox{\vrule width0pt height4.000000bp depth0pt\special{pdf:bcontent}\special{pdf:code q 0 g 0.500000 4.000000 m 0.000000 3.500000 l 0.000000 3.250000 l 0.750000 2.500000 l 0.750000 2.250000 l 1.000000 2.000000 l 1.000000 1.750000 l 0.750000 1.500000 l 0.250000 1.500000 l 0.000000 1.250000 l 0.000000 0.500000 l 0.250000 0.250000 l 1.250000 0.250000 l 1.500000 0.000000 l 2.000000 0.000000 l 2.250000 0.250000 l 2.500000 0.250000 l 3.250000 1.000000 l 3.250000 1.250000 l 2.500000 2.000000 l 2.500000 2.500000 l 3.000000 3.000000 l 3.000000 3.250000 l 2.250000 4.000000 l 1.750000 4.000000 l 1.500000 3.750000 l 1.250000 3.750000 l 1.000000 4.000000 l h f* Q}\special{pdf:econtent}\kern3.500000bp}}

\begin{center}\small Göttingische gelehrte Anzeigen. Stück 54. Seiten 531..536. 5. April 1817.\\[3pt]\rule{.72\linewidth}{.3pt}\end{center}

\emph{Connaissance des tems, ou des mouvemens célestes, à l'usage des astronomes et des navigateurs, pour l'an \textup{1818}; publiée par le bureau des longitudes.} Paris 1815. Bei der Witwe Courcier. 412 Seiten im Oktavformat.

\begingroup\brokenpenalty=10000\relax
Die Einrichtung des astronomischen Kalenders hat in diesem Jahrgang einige kleine Änderungen erfahren. Die Rektaszension des Mondes, die früher nur in Bogenminuten angegeben wurde, ist jetzt für Mittag und Mitternacht auf Sekunden berechnet. Die Deklinationen dagegen, die früher alle sechs Stunden angegeben wurden, sind jetzt nur für Mittag und Mitternacht in Minuten angegeben. Diese Änderung muss denjenigen Astronomen willkommen sein, die den Unterschied ihrer geografischen Längen durch korrespondierende Beobachtungen von Mondkulminationen bestimmen wollen. Ferner sind die einzelnen Phänomene, die früher für jeden Monat gesondert aufgeführt wurden, jetzt zweckmäßiger am Ende des Kalenders zusammengestellt. Der dadurch gewonnene Platz ist für die Angabe der heliozentrischen Längen und Breiten der Planeten und ihrer Rektaszensionen in Zeitmaß verwendet worden. Es ist nur zu bedauern, dass man die Rektaszensionen bloß auf Minuten (also auf Viertelgrade im Winkelmaß) angegeben hat. Um die Stellung der Planeten gegenüber benachbarten Fixsternen im Voraus beurteilen zu können, ist es sehr oft wünschenswert, die Position der Planeten auf die Bogenminute genau zu kennen. Die Angabe des Abstands der Planeten von der Erde wäre bei weitem wichtiger als die Angabe der heliozentrischen Positionen, die in der Ephemeride wenig oder gar keinen Nutzen haben. Es ist sehr zu bedauern, dass keine einzige unserer astronomischen Ephemeriden dieses Element mitteilt, obwohl der Abstand von der Erde für Aberration, Parallaxe und die Reduktion der gemessenen Durchmesser gleichermaßen wichtig ist.

Die beigefügten Aufsätze beginnen mit einigen Auszügen, die Herr \textsc{Delambre} aus verschiedenen englischen Schriften zusammengestellt hat. Da diese auf dem Festland noch immer nur in wenige Hände gelangen, werden die Auszüge manchen Lesern willkommen sein. Auch die beigefügten Anmerkungen eines so erfahrenen Praktikers liest man gern, wenn er sich auch häufig wiederholt. Mit ermüdender und unnötiger Weitschweifigkeit ist dagegen der folgende Aufsatz über eine Aufgabe von \textsc{Regiomontan} geschrieben. Die Aufgabe besteht darin, die Umstände anzugeben, unter denen sich zwei Punkte der Ekliptik in der Länge um ebenso viel unterscheiden wie in der Rektaszension. Diese unbedeutende und \SourceLeaf{601}leichte Aufgabe hätte in der Connoissance des tems keine 10 Seiten verdient. --- Neue Tafeln für die Aberration der Planeten in Länge und Breite von \textsc{Puissant}. Diese Tafeln sind so zweckmäßig eingerichtet, wie es die Natur des Gegenstands erlaubt, und können mit Nutzen angewandt werden, wenn der Abstand des Planeten von der Erde völlig unbekannt ist. Zu den Formeln, nach denen die Tafeln berechnet sind, bemerkt der Verfasser, Ces formules sont \emph{complètes}; il est remarquable en outre, qu'elles sont exactes, \emph{aux quantités près du troisième ordre}; car les termes en $\mathrm{e}^2$ se detruisent mutuellement, comme je m'en suis assuré. Vermutlich soll diese etwas dunkel formulierte Stelle bedeuten, dass der Verfasser die Aberration nach Potenzen der Exzentrizität entwickelt hat und dass sowohl die von der Exzentrizität unabhängigen Glieder als auch die Glieder erster Ordnung in der Exzentrizität vollständig sind. Wenn diese Auslegung zutrifft (ohne sie scheinen die unterstrichenen Ausdrücke im Widerspruch zu stehen), scheint dem Verfasser entgangen zu sein, dass seine Formeln tatsächlich über die Größen dritter Ordnung in der Exzentrizität hinaus genau und \emph{absolut} vollständig sind. --- Vom selben Verfasser eine Anmerkung zu einem Aufsatz von \textsc{Lagrange} über die Parallaxenrechnung bei Finsternissen. Herr \textsc{Puissant} führt Hilfsmittel ein, um den von \textsc{Lagrange} für den scheinbaren Abstand der Mittelpunkte von Sonne und Mond angegebenen Formeln eine für die numerische Rechnung geeignetere Gestalt zu geben. Bequemer ist hier jedoch die Verwendung der Tafel, aus der man unmittelbar die Logarithmen der Summen und Differenzen von Größen findet, die direkt durch ihre Logarithmen gegeben sind. --- Die neue allgemeine Tafel für die parabolische Bewegung der Kometen von \textsc{Burckhardt} verdient den Vorzug vor der gewöhnlichen Tafel und sogar vor der von \textsc{Barker}. Sie unterscheidet sich von der gewöhnlichen Tafel dadurch, dass der Logarithmus der seit dem Durchgang durch das Perihel verstrichenen Zeit an die Stelle dieser Zeit als Argument tritt. Gegenüber der gewöhnlichen Tafel erspart sie das Nachschlagen eines Logarithmus. Denselben Vorteil hat sie gegenüber der Tafel von \textsc{Barker} bis zu $45^\circ$ wahrer Anomalie, und auch das Interpolieren ist in der Tafel von \textsc{Burckhardt} meist etwas bequemer. --- Neue Bestimmung der Bahn der Vesta von \textsc{Daussy}. Herr \textsc{Daussy} hat sich große Verdienste erworben, indem er die Störungen der Bewegung der Vesta durch Jupiter, Saturn und Mars nach der Methode von \textsc{Laplace} berechnet hat. Bei der mäßigen Exzentrizität und Neigung dieses Planeten kann die Methode als ausreichend gelten. Die ersten sieben Oppositionen zeigen tatsächlich eine sehr befriedigende Übereinstimmung. --- Über einen neuen Apparat zum Vergleich von Längenmaßen, von \textsc{Prony}. Das wichtigste Stück dieses Apparats ist ein auf Glas in 100 Teile unterteilter Millimeter. Mehrere französische Instrumentenmacher stellen diese Teilung mit wahrhaft bewundernswerter Feinheit und Genauigkeit her. Die zu vergleichenden Maße müssen so zusammengesetzt werden, dass sich die Zusammensetzungen um höchstens ein paar Millimeter unterscheiden. Nur diese Differenz bestimmt der Apparat, dessen nähere Beschreibung hier zu weit führen würde. --- Über Ebbe und Flut des Meeres, von \textsc{Laplace}. Dieser kurze Aufsatz ist die Einleitung zu neuen, interessanten Untersuchungen, die dieser große Geometer über die in Brest angestellten Beobachtungen von Ebbe und Flut durchgeführt hat. Er muss sehr neugierig auf die Einzelheiten dieser Untersuchungen machen. Besonders bemerkenswert sind die daraus abgeleiteten Bestimmungen der Mondmasse $(= \frac{1}{68.7})$ und der Nutationskonstante $(= 9''65)$. Dies gilt umso mehr, als die von Herrn \textsc{von Lindenau} mit großer Sorgfalt untersuchten Beobachtungen des Polarsterns einen beträchtlich kleineren Wert für die Nutationskonstante ergeben haben. --- Die beiden folgenden Aufsätze desselben Verfassers über die Anwendung der Wahrscheinlichkeitsrechnung auf die Naturwissenschaft befassen sich mit der Theorie der Genauigkeit der nach der Methode der kleinsten Quadrate gefundenen Ergebnisse und mit der Bestimmung der Genauigkeit der Beobachtungen selbst. Die Ergebnisse dieser Untersuchung stimmen im Grunde ganz mit dem überein, was \textsc{Gauß} in der Theoria Motus \SourceLeaf{602}Corporum Coelestium und in einem Aufsatz in der Zeitschrift für Astronomie entwickelt hat, obwohl \textsc{Laplace} den Gegenstand aus einem etwas anderen Blickwinkel betrachtet hat. Zugleich werden hier mehrere interessante Anwendungen dieser Theorie angegeben. Eine davon ist die Bestimmung der Jupitermasse aus den Störungen, die Saturn durch Jupiters Einwirkung erfährt. Das aus dieser Quelle von \textsc{Bouvard} abgeleitete Ergebnis ist $\frac{1}{1070.5}$, und der wahrscheinliche Fehler dieser Bestimmung ergibt sich $= \pm \frac{1}{72\ObservedJupiterErrorMark}$ dieses Wertes. Das heißt, es ist ebenso wahrscheinlich, dass der wahre Wert zwischen $\frac{1}{1069}$ und $\frac{1}{1072}$ liegt, wie dass er außerhalb dieser Grenzen liegt. \textsc{Laplace} hat es vorgezogen, die außerordentlich geringe Wahrscheinlichkeit für eine Überschreitung etwas weiterer Grenzen anzugeben. Man kann 1 gegen fast eine Million wetten, dass der Fehler höchstens $\frac{1}{100}$ des Ganzen beträgt. Nun ergibt sich der höchst bemerkenswerte Umstand, dass die von \textsc{Gauß} aus seiner Theorie der Störungen der Pallas durch Jupiter abgeleitete Jupitermasse beträchtlich von jener Bestimmung abweicht und weit außerhalb jener Grenzen liegt. Die Hauptsache ist, dass die auf diese Bestimmung angewandte Wahrscheinlichkeitstheorie ihr eine beinahe ebenso geringe Unsicherheit zuschreibt. Darüber darf man sich nicht wundern, da bei den Bewegungen der Pallas der viel stärkere Einfluss Jupiters die noch geringere Zahl von Beobachtungen schon jetzt ausgleicht. Wie soll man nun diesen Widerspruch auflösen, und welches ist die wahre Masse Jupiters? Die fortgesetzten Beobachtungen der Pallas und die Untersuchung der Störungen der Juno werden darüber künftig weiteren Aufschluss geben. Es scheint aber berechtigt, daran zu zweifeln, ob die Methode von \textsc{Laplace} zur Berechnung der Störungen bei Anwendung auf Saturn ganz so genaue Ergebnisse liefert, wie sie für eine so heikle Untersuchung erforderlich sind. --- Über die Formeln von \textsc{Nonius} für die Dämmerung, von \textsc{Delambre}. Das Verdienst von \textsc{Nonius} um die Aufgaben, die sich auf die Dauer der Dämmerung beziehen, wurde bisher nicht genügend gewürdigt. Herr \textsc{Delambre} zeigt in dieser sehr ausführlichen Abhandlung, dass \textsc{Nonius} tatsächlich alle einschlägigen Aufgaben sehr gut und zum Teil besser als seine Nachfolger gelöst hat. Wir haben diese Abhandlung als eine Probe eines großen Werks anzusehen, das Herr \textsc{Delambre} als Geschichte der Astronomie herausgeben will. Darin sollen alle bekannten Bücher, wenigstens diejenigen, die etwas Nützliches oder Bemerkenswertes enthalten, in Auszügen und mit Erläuterungen, gelegentlich als Neudruck, in jedem Fall aber wenigstens in neuer Gestalt vorgelegt werden. --- Der letzte Aufsatz dieses Jahrgangs enthält die Elemente und Störungsgleichungen, nach denen die neuen Saturntafeln von \textsc{Bouvard} berechnet sind. Er enthält auch den Vergleich dieser Tafeln mit allen seit 1747 beobachteten Oppositionen und Quadraturen dieses Planeten, insgesamt 130. Wir hätten sehr gewünscht, dass alle diesen Tafeln zugrunde liegenden beobachteten Positionen mit abgedruckt worden wären. So gut die Tafeln auch die Beobachtungen darstellen, früher oder später wird man doch nach weiteren Verbesserungen suchen müssen, wobei man auf alle jene Oppositionen zurückkommen muss.
\par\clearpage\endgroup
